\documentclass{standalone}
\usepackage[utf8]{inputenc}
\usepackage[T1]{fontenc}
\usepackage{geometry}
\usepackage{graphicx}
\usepackage{float}
\usepackage{hyperref}
\usepackage[english,indonesian]{babel}

\title{Optimasi Pemrosesan Template LaTeX Menggunakan Pendekatan Berbasis AI}
\author{Academic Support Agent}
\date{\today}

\begin{document}
\maketitle

\begin{abstract}

\end{abstract}

\chapter{PENDAHULUAN}


\chapter{TINJAUAN PUSTAKA}


% Original content (commented out):
% \documentclass[svgnames]{standalone}
\usepackage[spanish]{babel}
\usepackage[T1]{fontenc}
\usepackage[utf8]{inputenc}
\usepackage{pgfplots, helvet, verbatim}
\renewcommand{\familydefault}{\sfdefault}
\pgfplotsset{width=10cm, compat=newest}
\usetikzlibrary{backgrounds}

\begin{comment}
Este ejemplo es la segunda versión del ejemplo "Espectro de Banda Base OFDM Multiportadora Logarítmica" https://www.overleaf.com/latex/examples/espectro-de-banda-base-ofdm-multiportadora-logaritmica/ccfysnxdbtsr. A diferencia de dicho ejemplo, en este caso se usó la opción gnuplot [raw gnuplot] de la orden \addplot. Al comparar ambas gráficas, se observan algunas diferencias, probablemente por la forma en que se calculan las coordenadas en el paquete PGFPLOTS y la opción GNUPLOT que usa órdenes del paquete homónimo.

Para el trazado de las 3 curvas uso un bucle \foreach y dentro del bucle, la orden \addplot+ con las opciones adecuadas para que las curvas aparezcan en colores diferentes. Las secciones comentadas explican la codificación de este gráfico.
\end{comment}
% Lista de colores en ciclo
\pgfplotscreateplotcyclelist{croma}{Goldenrod, Blue, Red}

% Sección colores de etiquetas
\newcommand{\titulo}{RoyalBlue} % Variable del color de la etiqueta del título.
\newcommand{\etiq}{VioletRed} % Variable del color de las etiquetas de los ejes.
\newcommand{\fondo}{LightGrey!15} % Variable del color del fondo del gráfico.

\begin{document}
\begin{tikzpicture}
\begin{axis}[
    legend style={fill=\fondo, font=\scriptsize,anchor=north east}, % Dibuja la caja de leyendas y aplica fondo
	axis background/.style={fill=\fondo},
	title style={text width=10cm},
	title=\centering \textbf{\color{\titulo} Espectro de Banda Base OFDM\\ Multiportadora Logarítmica 2\\},
	cycle list name = croma,
	xlabel=\color{\etiq} $\frac{f}{N*Rs}$,
	ylabel=\color{\etiq} $S(f)$,
	minor tick num=9,
	grid=major,
	ymin=-50
	]
	\foreach \i in {0, ...,2}
	\addplot +[no marks, line width=1, samples=100, smooth] gnuplot [raw gnuplot]{
	set xr[-0.1:2]; % Fija el rango del eje x
	f(x) = (sin(pi*x)/(pi*x))^2; % Definición de la función
	g0(x) = sum[n=0:3] f(4*x-n);
    g1(x) = sum[n=0:15] f(16*x-n);
    g2(x) = sum[n=0:63] f(64*x-n);
    plot 10*log10(g\i(x));
    };
	% Leyenda
	\legend{$N=4$, $N=16$, $N=64$}
\end{axis}
\end{tikzpicture}
\end{document}

\end{document}